\section{Study 1: Methods and Results}
\label{sec:study1}

\subsection{What we manipulate}

We ask one primary question. When a tutoring scaffold starts to fall apart
over multiple turns, does a second-pass oversight step hold the structure
better than putting the same rules in the prompt alone?

To answer that, we only change how ADHD-supportive structure is enforced.
Everything else stays fixed: the same chat model
(\texttt{google:gemini-2.5-flash} for the reported runs), the same retrieval,
tools, and temperature, and the same synthetic user turns.

That gives three arms (Figure~\ref{fig:threearms}):
\begin{enumerate}
  \item \textbf{Baseline:} ordinary tutoring; Assist off.
  \item \textbf{Assist, prompt-only:} ADHD Assist policy in the teacher system
    prompt; no second pass.
  \item \textbf{Assist + oversight:} same policy, plus a Dean that reads the
    full draft and revises when the constitution fails.
\end{enumerate}

The independent variable is not a smarter model or better retrieval. It is
whether structure is requested in the prompt, or checked after generation.

\subsection{What we measure}

We score each assistant turn for structural adherence with an automatic
checker.

\paragraph{Strict checklist.}
The reply must include a summary-first marker, end with a continue invite, and
stay under the word cap. Harsh on short redirects that correctly refuse a full
teaching template.

\paragraph{Turn-aware checklist (primary for Assist comparisons).}
Same idea, but a redirect may use the short boundary template instead of the
full tutoring scaffold. That is the fairer primary score when comparing
prompt-only to oversight.

\paragraph{Interface labels.}
The live UI may show TLDR and Continue where the system stores and scores Top
summary and Next?. Study~1 scores the internal anchors.

We also keep a late-turn aggregate (second-and-later turns in a scenario, and
the late tranche of the long probe), because that is where prompt drift shows
up if it is going to. The headline result is the automated pass rates on the
three-arm freeze.

\subsection{Probes}

All Study~1 inputs are synthetic multi-turn tutoring scenarios. No learner
accounts. No personal data. We average five independent repeats per arm. The
freeze includes a multi-turn topic interruption (core drift), continue-a-plan
after a break (re-orientation), and shorter and longer variants in the same
suite so we do not overfit to two scripts.

\subsection{How we ran the freeze}

Each arm ran as its own server configuration (oversight is a process-level
flag). We stored run metadata (git SHA, model id, arm label) with turn-level
results. For each arm we report the mean pass rate across five repeats, with
the observed range. These are measured compliance rates under a fixed harness,
not population estimates with confidence intervals.

\subsection{What Study 1 deliberately is not}

Study~1 does not measure cognitive load, learning gain, or preference; those
belong to the human protocol (\S\ref{sec:feasibility}). Study~1 asks whether
the scaffold stays on the page when we ask the system to keep it there.

\subsection{Results}
\label{sec:study1results}

Table~\ref{tab:freeze} reports mean structural pass rates over five repeats on
the frozen multi-turn suite (\texttt{google:gemini-2.5-flash}).

\begin{table}[t]
\caption{Structural adherence on the multi-turn probe suite (mean; ranges in
parentheses). Turn-aware scoring allows correct redirects to use a short
boundary template; strict scoring demands full tutoring markers on every
turn.}
\label{tab:freeze}
\footnotesize
\centering
\setlength{\tabcolsep}{4pt}
\begin{tabular}{lccc}
\toprule
Metric & Baseline & Prompt-only & Oversight \\
\midrule
Overall (turn-aware)
  & 0\% (0--0\%)
  & \textbf{76\%} (50--93\%)
  & \textbf{80\%} (71--86\%) \\
Late-turn (turn-aware)
  & 0\% (0--0\%)
  & \textbf{86\%} (71--100\%)
  & \textbf{89\%} (86--100\%) \\
Overall (strict)
  & 0\% (0--0\%)
  & 67\% (50--79\%)
  & 71\% (64--79\%) \\
Late-turn (strict)
  & 0\% (0--0\%)
  & 77\% (71--86\%)
  & 80\% (71--86\%) \\
\bottomrule
\end{tabular}
\end{table}

Baseline is \textbf{0\% (0--0\%)} on both checklists across five repeats:
unconstrained replies still answer the probes, but they never emit the
required ADHD-supportive shape. That zero is not a straw man: it is what the
deployed default tutor produces under the same probes, so the informative
contrast for this paper is prompt-only vs oversight. Prompt-only Assist
produces the large jump
(turn-aware \textbf{76\%} overall~/~\textbf{86\%} late; strict \textbf{67\%}~/
\textbf{77\%}). Oversight adds a further lift (turn-aware \textbf{80\%}~/
\textbf{89\%} late; strict \textbf{71\%}~/~\textbf{80\%}), about four points
in aggregate, with larger gains on residual hard turns (plan resume and
post-interrupt failures that the prompt arm still misses). Strict
rates run lower than turn-aware because intentional redirects fail a
full-marker checklist by design; we claim scaffold adherence here, not content
invariance under Dean rewrite.
